\documentclass[a4paper,10pt]{article}
%\usepackage[ngerman]{babel}
%\usepackage[top=3cm,bottom=3cm,left=3cm,right=3cm]{geometry} %NICHT IN KEMIE2


\usepackage{microtype} %schöneres typesetting  mit [stretch=10] für nur 1% anstelle von 2%
\usepackage{lmodern} %Latin Modern FONT
\usepackage{pifont} %schönere Font

\usepackage{chemfig}
\usepackage[version=4]{mhchem}
\usepackage{amsmath}
\usepackage{nicefrac}
\usepackage{amssymb}
%\usepackage{mathtools}

\usepackage{titlesec} %Um Section, etc bearbeiten zu können
\usepackage{fancyhdr}
\usepackage{setspace}
\usepackage{enumitem} %enumerate
\usepackage{arydshln} %für dashline
\usepackage{multicol}
\usepackage{lipsum}
\usepackage[labelfont=bf]{caption}

\usepackage[many]{tcolorbox}

%\usepackage{pgfornament} %Scheene Ornamente
\usepackage{witharrows} %Arrows in align-Umgebung
%\usepackage{awesomebox} %Notizboxen, etc
%\usepackage{textpos} %Für zum Beispiel vertikale Texte am Rand


%\usepackage{wrapfig} %Text um Bilder
%\usepackage{varwidth} % Alternative für engere minipage


\usepackage{hyperref}


\tcbuselibrary{theorems}

%\titlespacing*{\section}{0pt}{10mm}{3mm}
%\titlespacing*{\subsection}{0pt}{6mm}{3mm}
%\titlespacing*{\subsubsection}{0pt}{8mm}{3mm}


\renewcommand{\ttdefault}{lmtt} %Schriftart wird geändert zu Latin Modern Typewriter


\hypersetup{hidelinks,
			colorlinks=true,
			linkcolor=black,
			citecolor=miudiutex,
			urlcolor=miugray2
			}

\setlength\fboxsep{0pt}				
\setlength\columnsep{20pt}
\setlength{\parindent}{0pt}
\setcounter{section}{0}






%\definecolor{miudiutex}{HTML}{2f3d39}
\definecolor{miudiutex}{HTML}{1d2926}
\definecolor{miugray}{HTML}{b4cccf}
%\definecolor{miugray2}{HTML}{4d7365}
\definecolor{miugray2}{HTML}{6F8C81}
\definecolor{beige}{HTML}{F4F8D3}
\definecolor{white2}{HTML}{f5f7f8}
\definecolor{red}{HTML}{cf1f5c}
\definecolor{red2}{HTML}{9C191B}
\definecolor{green}{HTML}{00A36C}
\definecolor{delphinidin2}{HTML}{315794}
\definecolor{delphinidin}{HTML}{8B9EBA}
\definecolor{uniblau}{HTML}{004e9f}
\definecolor{unigelb}{HTML}{fcba00}
\definecolor{unigrau}{HTML}{909085}
\definecolor{pelargonidin}{HTML}{b33807}
\definecolor{pelargonidin2}{HTML}{FFDC96}
\definecolor{cyanidin}{HTML}{ba0746}
\definecolor{paeonidin}{HTML}{d15ca6}
\definecolor{petunidin}{HTML}{B88495}
\definecolor{petunitext}{HTML}{963354}
\definecolor{malvidin}{HTML}{8a2d8c}
\definecolor{mytheorembg}{HTML}{F2F2F9}
\definecolor{mytheoremfr}{HTML}{00007B}
\definecolor{uniblau2}{HTML}{27548A}
\definecolor{unigelb2}{HTML}{DDA853}
\definecolor{unidunkelblau2}{HTML}{183B4E}
\definecolor{unibeige}{HTML}{F5EEDC}


\tikzset{>=Latex}

\raggedcolumns %wichtig für Multicolumns
\color{miudiutex}



\usetikzlibrary{chains,
                positioning,
                shapes.multipart,
                }


\setchemfig
{%	
	bond offset=2pt,
	atom sep=15pt, 
	cram width=3pt,
	cram dash width=0.8pt,
	cram dash sep=1.2pt,
	double bond sep=2.5pt, 
	bond style={line width=1.1pt,cap=round},
	baseline=(current bounding box.center),
	bond join=true
}
\setcharge{extra sep=3pt}




\usepackage[ngerman]{babel}
\usepackage[top=3cm,bottom=3cm,left=3cm,right=3cm]{geometry} %NICHT IN KEMIE2


%\usepackage{microtype} %schöneres typesetting  mit [stretch=10] für nur 1% anstelle von 2%
%\usepackage{lmodern} %Latin Modern FONT
%\usepackage{pifont} %schönere Font



\usepackage{amsmath}
\usepackage{nicefrac}
\usepackage{amssymb}
%\usepackage{mathtools}

\usepackage{titlesec} %Um Section, etc bearbeiten zu können
\usepackage{setspace}
\usepackage{enumitem} %enumerate
\usepackage{multicol}
\usepackage[many]{tcolorbox}
\usepackage{hyperref}
\usepackage{arydshln}

%\renewcommand{\ttdefault}{lmtt} %Schriftart wird geändert zu Latin Modern Typewriter


\hypersetup{hidelinks,
			colorlinks=true,
			linkcolor=black,
			citecolor=miudiutex,
			urlcolor=miugray2
			}

\setlength{\parindent}{0pt}
\setcounter{section}{0}


\titlespacing*{\section}{0pt}{20mm}{6mm}
\titlespacing*{\subsection}{0pt}{6mm}{3mm}
\titlespacing*{\subsubsection}{0pt}{3mm}{0mm}



%\definecolor{miudiutex}{HTML}{2f3d39}
\definecolor{miudiutex}{HTML}{1d2926}
\definecolor{miugray}{HTML}{b4cccf}
\definecolor{miugray2}{HTML}{4d7365}
\definecolor{red}{HTML}{cf1f5c}
\definecolor{green}{HTML}{00A36C}
\definecolor{delphinidin}{HTML}{315794}
\definecolor{pelargonidin}{HTML}{b33807}
\definecolor{cyanidin}{HTML}{ba0746}
\definecolor{petunidin}{HTML}{b56fed}



\raggedcolumns %wichtig für Multicolumns


%================================
% SYMBOLS
%================================

\newcommand{\RA}{$\rightarrow$~}
\newcommand{\RL}{$\rightleftharpoons$~}
\newcommand{\HARP}{\ding{226}~}
%$\xrightarrow{2}$ %Pfeil mit 2 drüber


%================================
% SHORT COMMANDS (BOXES ; ETC)
%================================


\newcommand{\notiz}[1]{\begin{notizz}{}{}\begin{center}
			{\color{miugray2!50!miudiutex}#1} \end{center}\end{notizz}}
\newcommand{\zit}[2]{\begin{zitat*}{#1}{}\small #2 \end{zitat*}}

\newcommand{\frage}[2]{\begin{qst*}{#1}{}\small #2 \end{qst*}}

\newcommand{\unten}{\end{center} \tcblower \begin{center}}


\newcommand*\circled[1]{\tikz[baseline=(char.base)]{
            \node[shape=circle,draw,inner sep=0.5pt,miudiutex] (char) {\tiny #1};}}
            
\newcommand{\miusquare}{{\color{miugray2!50}\scriptsize $\blacksquare$}~}
\newcommand{\miusquarp}{{\color{petunidin!50}\scriptsize $\blacktriangleright$}~}
\newcommand{\niu}{\item[\miusquare]}
\newcommand{\nip}{\item[\miusquarp]}

%%%%% ABSAETZE

\newcommand{\pps}{\par \vspace*{1mm}}
\newcommand{\pp}{\par \vspace*{3mm}}
\newcommand{\ppbig}{\par \vspace*{8mm}}
\newcommand{\pphuge}{\par \vspace*{10mm}}

%================================
% FRAGE BOX
%================================

\tcbuselibrary{theorems,skins,hooks}
\newtcbtheorem{qst}{Frage:}
{%
	enhanced,
	breakable,
	colback = gray!15!white,
	frame hidden,
	boxrule = 0sp,
	borderline west = {2.5pt}{0pt}{red!70},
	sharp corners,
	detach title,
	before upper = \tcbtitle\par\smallskip,
	coltitle = black,
	fonttitle = \bfseries\sffamily,
	description font = \mdseries,
	separator sign none,
	segmentation style={solid, gray!50!black},
}
{th}




%================================
% ZITAT BOX
%================================

\tcbuselibrary{theorems,skins,hooks}
\newtcbtheorem{zitat}{Zitat}
{%
	enhanced,
	breakable,
	colback = gray!15!white,
	frame hidden,
	boxrule = 0sp,
	borderline west = {2.5pt}{0pt}{black!70},
	sharp corners,
	detach title,
%	before upper = \tcbtitle\par\smallskip,
	coltitle = gray!50!black,
	fonttitle = \bfseries\sffamily,
	description font = \mdseries,
%	separator sign none,
	segmentation style={solid, gray!50!black},
}
{th}



%================================
% SIMPLE SCHWARZER RAHMEN BOX
%================================


\newcommand{\boxfr}[1]{
	\begin{tcolorbox}[
	drop shadow southeast,
	enhanced,
	colframe=black!50,
	colback=white,
	boxrule = 1pt, 
%	toprule=0pt, bottomrule=0pt,rightrule=0pt,leftrule=0pt,
	boxsep=5pt,
	arc=1pt,
	]
	\begin{center}
	#1
	\end{center}
	\end{tcolorbox}
}

\usepackage[utf8]{inputenc}
\usepackage[T1]{fontenc}
\usepackage{lmodern}

\usepackage[
    activate={true,nocompatibility},
    final,
    tracking=true,
    kerning=true,
    spacing=true,
    factor=1100,
    stretch=30,
    shrink=20
]{microtype}
\usepackage{pdfpages}
\usepackage{awesomebox}
\usepackage{marginnote}
\tcbuselibrary{documentation}
\usepackage{sidenotes}
\usepackage{import}
\usepackage{xifthen}
\usepackage{transparent}
\usepackage[table]{xcolor}


\newcommand{\abbicon}{%
  \protect\tikz[baseline=0.2mm,scale=0.8]{%
    % Das blaue Quadrat mit abgerundeten Ecken
    \fill[uniblau, rounded corners=1.2pt] (0,0) rectangle (0.8em, 0.8em);
    % Der weiße Kreis in der Mitte
    \fill[white] (0.4em, 0.4em) circle (0.12em);
  }%
} 
\usepackage[labelfont={bf},font={color=miudiutex,small},figurename={\abbicon $~$Abbildung}]{caption}


\newcommand{\bckgrnd}{
\AddToHookNext{shipout/background}{
\begin{tikzpicture}[remember picture, overlay]
 \draw[draw=none,fill=uniblau!70, shift={(current page.north west)}] (-3,0) rectangle (23,-3.75);
 \draw[draw=none,fill=miugray, shift={(current page.north east)}] (-7,0) rectangle (3,-3.75);
\end{tikzpicture}
}
}



\newcommand{\incfig}[2]{%
\def\svgwidth{#2\columnwidth}
\import{C://LaTeX-Files/pngs/}{#1.pdf_tex}
}

\renewcommand{\textbf}[1]{\textsf{{\bfseries {#1}}}}



\titleformat{\section}[hang]
		{\LARGE \color{uniblau!75!black} \sffamily \bfseries}
		{\thesection}
		{6mm}
		{}
		[]
\titleformat{\subsection}[hang]
		{\large \sffamily \bfseries \color{black}}
		{\thesubsection}
		{4mm}
		{{\color{miugray}$\blacksquare ~$}}

\titleformat{\subsubsection}[block]
		{\normalsize \bfseries \sffamily \color{miudiutex}}%
		{\normalsize \color{black} \thesubsubsection}
		{2mm}
		{{\color{miugray}$\blacksquare ~$}}
		[ \color{black}]
		
\titleformat{\paragraph}
		{\normalsize \bfseries \sffamily \color{black}}
		{\footnotesize \theparagraph}
		{1mm}
		{}

\pagestyle{empty}
\titlespacing*{\section}{0pt}{20mm}{12mm}
\titlespacing{\section}{0pt}{20mm}{12mm}

\titlespacing*{\subsection}{0pt}{14mm}{4mm}
\titlespacing{\subsection}{0pt}{14mm}{4mm}

\titlespacing*{\subsubsection}{0pt}{6mm}{3mm}
\titlespacing{\subsubsection}{0pt}{6mm}{3mm}

\titlespacing*{\paragraph}{0pt}{6mm}{2mm}
\titlespacing{\paragraph}{0pt}{6mm}{2mm}


\let\oldmarginfigure\marginfigure
\let\endoldmarginfigure\endmarginfigure

\let\oldmarginfigure\marginfigure
\let\endoldmarginfigure\endmarginfigure

\renewenvironment{marginfigure}[1][0pt] % [1] steht für 1 Argument, [0pt] ist der Standardwert
  {\oldmarginfigure[#1]\captionsetup{labelfont={sf,bf,color=uniblau!75!black},font={color=miudiutex,small}}}
  {\endoldmarginfigure}
  
  
  
  
\renewcommand{\labelitemi}{\color{uniblau!70}$\bullet$}
\renewcommand{\labelitemii}{\color{uniblau!70}$\cdot$}
\newcommand{\plmtsquare}{{\color{uniblau!70}$\blacksquare~$}}
\newcommand{\splmt}[1]{{\color{uniblau!75!black} \sffamily #1}}
\newcommand{\fplmt}[1]{{\color{uniblau!75!black} \sffamily \bfseries #1}}

\newcommand{\antwort}[1]{
\begin{tcolorbox}[
	enhanced,
	enforce breakable,
	standard jigsaw,
	boxrule=0.7pt,
	colframe=black,
	sharp corners,
	boxsep=5pt,
	title=\bfseries \sffamily Antwort,
	opacityback=0,
	coltitle=miudiutex,
	attach boxed title to top text left={yshift=-\tcboxedtitleheight/2},
	boxed title style={colback=white,colframe=white},
	bottomrule at break=0pt,
	toprule at break=0pt,
	pad at break=16mm,
]
	\begin{center}
	#1
	\end{center}
	\end{tcolorbox}
}




\rowcolors{2}{miugray!50}{miugray!50}
\arrayrulecolor{white} % Weiße Gitterlinien
\setlength{\arrayrulewidth}{1.5pt} % Etwas dickere Linien
\renewcommand{\arraystretch}{1.4} % Mehr Zeilenhöhe für bessere Lesbarkeit}


%\setlist[itemize]{itemsep=0mm}
\setlist[enumerate]{itemsep=0mm}
%\usepackage[fleqn]{amsmath}
\usepackage{witharrows}
\usepackage{awesomebox}
\usepackage{amssymb}
\usepackage{siunitx}
\usepackage{setspace}
%\usepackage[top=1.2cm,bottom=1.2cm,left=2cm,right=2cm]{geometry}
\renewcommand{\textbf}[1]{\textsf{{\bfseries {#1}}}}
\setlength{\parindent}{0pt}
\setlength{\mathindent}{0cm}
\newcommand{\cm}{\text{cm}}
\newcommand{\m}{\text{m}}
\newcommand{\lite}{\text{l}}
\newcommand{\kg}{\text{kg}}
\newcommand{\km}{\text{km}}
\newcommand{\h}{\text{h}}
\newcommand{\s}{\text{s}}
\newcommand{\tonne}{\text{t}}

\newcommand\ddfrac[2]{\frac{\displaystyle #1}{\displaystyle #2}}

\pagestyle{empty}

\definecolor{miugray}{HTML}{b4cccf}
\definecolor{uniblau}{HTML}{004e9f}
\renewcommand{\plmtsquare}{{\color{miugray} $\blacksquare~$}}
\begin{document}

%\pagebreak
\newgeometry{top=3cm,bottom=3cm,left=3cm,right=3cm}
\vspace*{\fill}
\Huge \sffamily Hinweis: \normalsize
\antwort{
Was vorausgesetzt wird und deshalb nicht Bestandteil dieser Formelsammlung ist:
\begin{itemize}

\item den Umgebungsdruck kennen
\item die Einheit bar in die Einheit Pascal umrechnen können
\item Mischungsgleichung und Mischungskreuz kennen und anwenden können
\item Massenwirkungsgesetz aufstellen können und Definition der Säurekonstante
\item pH-Formel für schwache Säuren
\item die Definition des Octanol-Wasser-Koeffizienten kennen
\item Massenbilanzen aufstellen können
\item einfache Berechnungen durchführen können, z.B. Fläche eines Kreises, Volumen eines Zylinders, etc.
\end{itemize}

}
\vspace*{\fill}
\pagebreak



\newgeometry{top=1.2cm,bottom=1.2cm,left=2cm,right=2cm}
\AddToHookNext{shipout/background}{
\begin{tikzpicture}[remember picture, overlay]
   \draw[draw=uniblau!70,fill=uniblau!70, shift={(current page.north west)}] (-3,0) rectangle (23,-2.75);
   \draw[draw=miugray,fill=miugray, shift={(current page.north west)}] (-3,0) rectangle (5,-2.75);
%   \draw[petunitext,fill=petunitext, shift={(current page.north west)}] (15.2,-2) rectangle (15.4,-5);
%	\node(b) at (10.3,-8) {\includegraphics[width=21.3cm]{F://LaTeX-Files/pngs/tomat}};
%	\node[rotate=0,white!35, align=left,scale=2.5] (a)at(10,-16){\bfseries \sffamily  Zusammenfassung \\ \LARGE \bfseries \sffamily Lebensmittelsicherheit};
%	\node(c) at (16.3,-26) {\includegraphics[width=5cm]{F://LaTeX-Files/pngs/banana}};
\end{tikzpicture}
}


\begin{center}
\LARGE
\color{white}
\textbf{Formelsammlung} \par \vspace*{2mm}
\large \color{miugray} \textbf{Prozessbezogene Lebensmitteltechnologie}
\end{center}
\par 
\vspace*{12mm}

\begin{minipage}[t]{.55\textwidth}%
%\large

\raggedbottom

\plmtsquare \textbf{Flüssigkeiten \& Gase}\par \vspace*{-4mm}
%\begin{spacing}{1.5}
\begin{align*}
&\text{Druck:} ~ p = \frac{F}{A} \\
&\text{Kontinuitätsgleichung:} ~ A_1 \cdot v_1 = A_2 \cdot v_2 \\
&\text{Bernoulli:} 		~ p + \rho g h + \frac{1}{2} \rho v^2 = \text{const.} \\[-4pt]
&\quad \text{Staudruck (nach Bernoulli):} ~ p = \frac{1}{2} \rho v^2 \\[-4pt]
&\quad \text{Druck in Wassersäule:} ~ p = \rho gh \\
&\text{Toricelli (Ausfließgeschw.)}: ~ v = \sqrt{2gh} \\[-2pt]
&\qquad t = \sqrt{\frac{2}{g}} \cdot \frac{A_1}{A_2} \cdot (\sqrt{h_{\alpha}}-\sqrt{h_{\omega}})
\end{align*}

\vspace*{8mm}
\plmtsquare \textbf{Strömung mit Reibung}\par \vspace*{-4mm}
\begin{align*}
&\text{Volumenstrom:}	~ \dot{V} = \frac{V}{t} = A \cdot \overline{v}\\
&\text{Schubspannung:}	~ \tau = \eta \cdot \frac{d v}{d z} = \frac{F_{R}}{A} \\
&\text{Viskosität:}	~ \eta = \frac{\tau}{D} \\
&\text{Hagen - Poiseuille:} ~ \dot{V} = \frac{\pi \cdot r^4}{8\eta} \cdot \frac{d p}{d x}\\
%&\text{Leistung} ~ P = \frac{8\eta \cdot \Delta l}{\pi \cdot r^4} \cdot \dot{V}^2\\
&\text{Druckverlust durch Reibung, laminar}\\
& \qquad ~ \Delta p = \frac{8\eta \cdot \Delta l}{\pi \cdot R^4} \cdot \dot{V}\\
&\text{Geschwindigkeitsprofil laminare Strömung:}\\
&\qquad ~ v(r) = -\frac{R^2}{4\eta} \cdot \frac{dp}{dx} \cdot \left[ 1 - \left(\frac{r}{R}\right)^2 \right]\\
&\text{mittlere}\\[-10pt] &\text{Strömungsgeschwindigk.:} ~  c = \overline{v} = \frac{1}{2} \cdot v_{max}\\ 
&\text{Stokes'sches Gesetz:} ~ F_r = 6 \pi \eta v r \\[-4pt]
&\quad \text{Stationäre Sinkgeschwindigk.:} ~ v = \frac{2r^2g(\rho_1-\rho_2)}{9\eta} \\[5pt]
&\text{Gesamtwiderstand:} ~ F_{tot} = F_r + F_w\\[-4pt]
&\quad \text{Reibungswiderstand:}  ~ F_r = \eta \cdot A \cdot \frac{d v}{d z}\\[-2pt]
&\quad \text{Widerstandskraft:}  ~ F_w = c_w \cdot A \cdot \frac{1}{2} \rho v^2 \\
\\ 
\end{align*}
\vspace*{-12mm}

\plmtsquare \textbf{Rheologie }\par 
\resizebox{.7\columnwidth}{!}{
\begin{tikzpicture}
\node(a)at(0,0){
\begin{tikzpicture}
\node[align=center,scale=0.8](1)at(0.5,1.5){Newtonsch};
\draw (0,1)node[right]{$\tau$} -- (0,0) -- (1,0)node[right,scale=0.8]{$\dot{\gamma}$};
\draw[yshift=-1.5cm] (0,1)node[right]{$\eta$} -- (0,0) -- (1,0)node[right,scale=0.8]{$\dot{\gamma}$};
\draw[black!80] (0,0) -- (0.9,0.9);
\draw[black!80,yshift=-1.5cm] (0,0.5) -- (0.9,0.5);
\end{tikzpicture}
};
\node(b)at(2.125,0){
\begin{tikzpicture}
\node[align=center,scale=0.8](1)at(0.5,1.5){Scherverdünnend};
\draw (0,1)node[right]{$\tau$} -- (0,0) -- (1,0)node[right,scale=0.8]{$\dot{\gamma}$};
\draw[yshift=-1.5cm] (0,1)node[right]{$\eta$} -- (0,0) -- (1,0)node[right,scale=0.8]{$\dot{\gamma}$};
\draw[black!80] (0,0) to [bend left=30] (0.9,0.9);
\draw[black!80,yshift=-1.5cm] (0,0.9) to [bend right=45] (0.9,0.3);
\end{tikzpicture}
};
\node(c)at(4.5,0){
\begin{tikzpicture}
\node[align=center,scale=0.8](1)at(0.5,1.5){Scherverdickend};
\draw (0,1)node[right]{$\tau$} -- (0,0) -- (1,0)node[right,scale=0.8]{$\dot{\gamma}$};
\draw[yshift=-1.5cm] (0,1)node[right]{$\eta$} -- (0,0) -- (1,0)node[right,scale=0.8]{$\dot{\gamma}$};
\draw[black!80] (0,0) to [bend right=30] (0.9,0.9);
\draw[black!80,yshift=-1.5cm] (0,0.3) to [bend right=30] (0.9,0.9);
\end{tikzpicture}
};
\node[scale=0.7,rotate=90](fl)at(-1,0.55){Fließkurve};
\node[scale=0.7,rotate=90,align=left](vi)at(-1,-1){Viskositäts-\\kurve};
\end{tikzpicture}
}
\vfill
\end{minipage}%
\hfill
\begin{minipage}[t]{.35\textwidth}

\plmtsquare \textbf{Einheiten \& Konstanten}\par \vspace*{-4mm}

\begin{align*}
%\text{Joule [J]} ~ 	= \frac{\kg \cdot \m^2}{\s^2} \\ %AUSGEBLENDET
&\text{Newton [N]} ~	= \frac{\kg \cdot \m}{\s^2} \\
&\text{Pascal [Pa]} ~	= \frac{\kg}{\m \cdot \s^2} = \ddfrac{\text{N}}{\m^2} \\ 
&\text{Viskosität [$\eta$]} ~	= \text{Pa}\cdot \s = \frac{\text{N} \cdot \s}{\m^2} \\ 
&\text{Watt [W]} ~ = \frac{\kg \cdot \m^2}{\s^3} = \frac{J}{s}\\ \\ \\
&\text{Kritische Reynoldszahl:} \quad \text{Re}_{\text{krit}} = \frac{\rho d v}{\eta} = 2.300 \\[-4pt] 
&\text{Stefan-Boltzmann Konst.:} ~\sigma = 5,67 \cdot 10^{-8}~ \text{W}/ \text{m}^{2}\text{K}^{4}\\[1pt]
&\text{Allgemeine Gaskonstante:} ~R = 8,314 ~~\text{J/mol K} \\
&\text{Strahlungszahl schw. Körper:} ~C_S = 5,67~ \text{W}/ \text{m}^{2}\text{K}^{4}\\
&\text{Erdbeschleunigung:} \quad g = 9,81 ~\text{m/s}^2 \\[-1pt]
&\text{Dichte von Wasser:} \quad 1.000 ~\text{kg/m}^3 \\
\end{align*}


\vspace*{8mm}
\plmtsquare \textbf{Rohre \& Pumpen:}\par \vspace*{-4mm}
\begin{align*}
&\text{Stokes im Rohr:} ~ v = \frac{\Delta p}{4\eta \cdot l} \cdot (r_a^2-r_i^2) \\ 
& \hspace*{3cm} v_{max} = \frac{\Delta p \cdot r^2}{4\eta \cdot l \eta} = 2 \cdot \overline{v} \\
&\text{Turbulente Strömungsgeschwindigkeit:}\\ & \quad v = v_{max} \cdot (1-\frac{r_i}{r_a})^{1/7}\\[-8pt]
&\hspace*{0.8cm} \scriptstyle{v_{max} \approx 1,2 \cdot \overline{v}}\\
&\text{Volumenstrom:} \quad \dot{V} = \frac{\pi \cdot \Delta p \cdot r^4}{8 \cdot l \cdot \eta} = A_{quer} \cdot \overline{v} \\ 
& \text{Druckabfall im Rohr:} ~ \Delta p = \xi \cdot \frac{l}{d} \cdot \frac{1}{2}\rho\overline{v}^2\\
&\text{Rohrreibungszahl Xi:} ~ \xi_{lam} = \frac{64}{\text{Re}}\\
&\text{Hydraulischer Durchmesser:} ~ d_h = 4\cdot \frac{A}{U}\\
&\text{Nutzleistung Pumpe:} ~ P = \dot{V} \cdot \Delta p_{ges}  \\
\end{align*}




\vspace*{8mm}
\plmtsquare \textbf{Strahlung} \par \vspace*{-4mm}
\begin{align*}
&\text{Schwarzkörperstrahlung:} \quad P = \varepsilon \cdot \sigma \cdot A \cdot T^4 \\[-6pt]
&\quad \text{Vereinfacht:} \hspace*{2cm} \quad P = \varepsilon \cdot C_S \cdot A \cdot (\frac{T}{100})^4 \\
&\text{Elektromagn. Strahlung:} \quad E = h \cdot c \cdot \ell^{-1} = h \cdot f \\
%&\text{Plank'sches Gesetz:} \quad i = \frac{C_1}{\ell^5 \cdot [\text{exp}(C_2/\ell \cdot T)-1]}
\end{align*}
\end{minipage}%
\vfill
	
	
\pagebreak
\raggedbottom

\begin{minipage}[t]{.45\textwidth}
\vspace*{5.5mm}
\plmtsquare \textbf{Wärmelehre} \par \vspace*{-4mm}
\begin{align*}
&\text{Wärmeenergie:} \quad Q = c_m \cdot m \cdot \Delta T \\
&\text{Wärmestromdichte:} \quad \dot{q} = \frac{\dot{Q}}{A} \\
&\text{Konvektion:} \quad \dot{Q} = \alpha \cdot A \cdot \Delta T \\[6pt] 
&\text{Wärmeleitung nach Fourier:}\\
&\quad \text{1.Gesetz:} \quad \dot{Q} = \lambda \cdot A \cdot \frac{\Delta T}{x} \\
&\quad \text{2.Gesetz:} \quad \frac{dT}{dt} = \frac{\lambda}{\rho \cdot c_p} \cdot \frac{d^2T}{dx^2} = a \cdot \frac{d^2T}{dx^2} \\ \\
&\text{Abschätzung Wärmeleitfähigkeit von LM:}\\[4pt]
& \hspace*{2cm} \lambda = \sum \lambda_i \cdot X_Vi \\ \\
&\text{Stat. Wärmeleitung d. Wand mit i Schichten:}\\[4pt]
&\hspace*{2cm} \dot{Q} = \frac{A \cdot \Delta T_{ges}}{\sum \frac{l_i}{\lambda_i}} = \frac{A \cdot \Delta T_{ges}}{R_{ges}}\\ \\
&\text{Stat. Wärmeleitung d. Rohr mit n Schichten:}\\[4pt]
&\hspace*{2cm} \dot{Q} = \frac{2 \pi \cdot l \cdot \Delta T }{\sum \frac{\ln(\frac{r_{a_n}}{r_{i_n}})}{\lambda_n}} \\ \\
&\text{Fläche einer Rohrwand:}\\[4pt]
& \hspace*{2cm} \overline{A} = \frac{A_a - A_i}{\ln(\frac{A_a}{A_i})} = \frac{2 \pi \cdot l \cdot (r_a - r_i)}{\ln(\frac{A_a}{A_i})} \\
\end{align*}



\vspace*{8mm}
\plmtsquare \textbf{Pasteurisieren, Haltbarmachen} \par \vspace*{-4mm}
\begin{align*}
&\text{Arrhenius-Gleichung:} \quad k = A \cdot e^{-\frac{E_A}{R \cdot T}} \\
&\text{Q$_{10}$-Wert:} \quad Q_{10} = \frac{k_{T+10}}{k_T} \\
&\text{z-Wert:} \quad z = T_{10k} - T_k \\
&\text{Kinetik 1.Ordnung:} \quad \frac{dN}{dt} = -k \cdot N \\
& \hspace*{3.3cm} N(t) = N_0 \cdot e^{-kt} \\
&\text{Dezimalreduktionszeit:} \quad D_1 = D_2 \cdot 10^{\frac{T_2 - T_1}{z}} \\[2pt]
&\hspace*{3.3cm} D = \frac{t_2-t_1}{\log(\frac{N_{(t_1)}}{N_{(t_2)}})} \\
&\text{F-Wert:} \quad F_1 = F_2 \cdot 10^{\frac{T_2 - T_1}{z}} \\
&\text{Inaktivierungsszeit:} \quad F_T = \log(\frac{N_0}{N}) \cdot D_T \\
&\text{Risiko für Kontamination:} \quad \frac{1}{n} = \frac{N_0}{10^{\frac{F}{D}}} 
\end{align*}



\end{minipage}\hfill
\begin{minipage}[t]{.45\textwidth}
\vspace*{5.5mm}
\plmtsquare \textbf{Diffusion, Packstoffe} \par \vspace*{-4mm}
\begin{align*}
&\text{Säurekonstante:} \quad K_S = \frac{[H^+] \cdot [A^-]}{[HA]} \\ 
&\text{Fick'sches Gesetz:} \quad \dot{n} = D \cdot A \cdot \frac{\Delta c}{l} \\[5pt]
&\text{Henry Gesetz:} \quad c_i = \alpha \cdot p_i \\[3pt]
&\text{Permeation:} \quad J = P \cdot A \cdot \frac{\Delta p}{l}   
\end{align*}


\vspace*{6mm}
\plmtsquare \textbf{Verdampfung} \par \vspace*{-4mm}
\begin{align*}
&\text{Clausius-Clapeyron:}\\[4pt]
& \quad \ln\left(\frac{p_2}{p_1}\right) = \frac{\Delta h_v \cdot M_L}{R} \cdot \left(\frac{1}{T_1} - \frac{1}{T_2}\right) \\ 
& \hspace*{17mm} = \frac{\Delta h_{vM}}{R} \cdot \left(\frac{1}{T_1} - \frac{1}{T_2}\right) \\ \\
& \text{Siedepunkterhöhung durch gelöste Stoffe:} \\[4pt]
& \quad \Delta T \approx \frac{R \cdot T_L^2}{\Delta h_{vM}} \cdot x_i = \frac{R \cdot T_L^2}{M_L \cdot \Delta h_v} \cdot x_i \approx \frac{R \cdot T_L^2}{\rho_L \cdot \Delta h_v} \cdot c_i \\ \\
& \text{Siedepunkterhöhung durch hydrostatischen Druck:} \\[4pt]
& \quad \Delta T = \frac{\rho_{lg} \cdot g \cdot h \cdot R \cdot T_L^2}{\Delta h_{vM} \cdot p} = \frac{\rho_{lg} \cdot g \cdot h \cdot R \cdot T_L^2}{M_L \cdot \Delta h_v \cdot p}\\ \\
& \text{Siedepunkterhöhung durch Oberflächenspannung:} \\[4pt]
& \quad \Delta T = \frac{4\sigma \cdot R \cdot T_L^2}{d \cdot \Delta h_{vM} \cdot p} \\ \\
& \text{Benötigter Wärmestrom für kontinuierliche Verdampfung:} \\
& \quad \dot{Q} = \dot{m}_F \cdot c_{pF} \cdot (T_S - T_F) + \dot{m}_D \cdot \Delta h_v
\end{align*}


\vspace*{4mm}
\plmtsquare \textbf{Kennzahlen turbulente Strömung:} \par \vspace*{-4mm}
\begin{align*}
&\text{Reynoldszahl:} \quad Re = \frac{v \cdot l \cdot \rho}{\eta} = \frac{v \cdot l}{\color{red}{\nu}}
\\
&\text{Grashofzahl:} \quad Gr = - \frac{g \cdot l^3}{{\color{red}{\nu}}^2} \cdot \frac{\Delta \rho}{ \rho} \\
&\text{Nusseltzahl:} \quad Nu = \frac{\alpha \cdot l}{\lambda} \\
&\text{Prandtlzahl:} \quad Pr = \frac{{\color{red}{\nu}}}{a} = \frac{\eta \cdot c_p}{\lambda}
\end{align*}
\par


{\color{red} \textit{Kinematische Viskosität $\nu$ in Rot dargestellt.}}
\par \vspace*{4mm}
\textbf{Anmerkung:} Bei diesen Formeln steht $l$ für die charakteristische Länge des jeweiligen geometrischen Objekts. Für Rohre wäre $l$ dann beispielsweise der Innendurchmesser $d$
\par \vspace*{8mm}


\vfill
\end{minipage}	

%\pagebreak
%\newgeometry{top=3cm,bottom=3cm,left=3cm,right=3cm}
%
%\antwort{
%Was vorausgesetzt wird und deshalb nicht Bestandteil dieser Formelsammlung ist:
%\begin{itemize}
%
%\item den Umgebungsdruck kennen
%\item die Einheit bar in die Einheit Pascal umrechnen können
%\item Mischungsgleichung und Mischungskreuz kennen und anwenden können
%\item Massenwirkungsgesetz aufstellen können und Definition der Säurekonstante
%\item pH-Formel für schwache Säuren
%\item die Definition des Octanol-Wasser-Koeffizienten kennen
%\item Massenbilanzen aufstellen können
%\item einfache Berechnungen durchführen können, z.B. Fläche eines Kreises, Volumen eines Zylinders, etc.
%\end{itemize}
%
%}


\end{document}
